\documentclass[10pt,reqno]{amsart}

\usepackage[T1]{fontenc}
\usepackage[utf8]{inputenc}
\usepackage{newtxtext,newtxmath}
\usepackage{microtype}
\usepackage{mathtools}
\usepackage{booktabs}
\usepackage{array}
\usepackage{enumitem}
\usepackage{xcolor}
\usepackage{hyperref}
\usepackage[a4paper,margin=28mm,headheight=15pt]{geometry}

\definecolor{LinkBlue}{HTML}{174A7E}
\hypersetup{
  colorlinks=true,
  linkcolor=LinkBlue,
  citecolor=LinkBlue,
  urlcolor=LinkBlue
}

\setlength{\parskip}{2pt}
\setlist[itemize]{leftmargin=6mm,itemsep=1.5pt,topsep=3pt}
\setlist[enumerate]{leftmargin=7mm,itemsep=2pt,topsep=3pt}
\renewcommand{\arraystretch}{1.15}
\emergencystretch=2em

\newtheorem{theorem}{Theorem}[section]
\newtheorem{proposition}[theorem]{Proposition}
\newtheorem{lemma}[theorem]{Lemma}
\newtheorem{corollary}[theorem]{Corollary}
\theoremstyle{definition}
\newtheorem{definition}[theorem]{Definition}
\newtheorem{remark}[theorem]{Remark}

\title[A Finite $K(G,1)$ Without Baumslag--Solitar Subgroups Can Be Non-Hyperbolic]
{A Finite $K(G,1)$ Without Baumslag--Solitar Subgroups Can Be Non-Hyperbolic}
\author{Argus Open Research Project}
\date{August 30, 2026}
\subjclass[2020]{Primary 20F67; Secondary 57M07}
\begin{document}
\begin{abstract}
Question~1.1 of Bestvina's problem list in geometric group theory asks whether every group $H$ with a finite classifying space $K(H,1)$ (or every finitely presented subgroup $H$ of a word-hyperbolic group $G$) that contains no Baumslag--Solitar subgroup $\mathrm{BS}(m,n)$ must be word-hyperbolic. We document and verify the simultaneous counterexample of Italiano, Martelli, and Migliorini (2021/2023): the kernel $H=\ker(f:G\to\mathbb{Z})$ of a fibration of a hyperbolic five-manifold over $S^1$ with a finite aspherical four-manifold fiber $F$.
\end{abstract}
\maketitle
\section{Main Theorem and Scope Boundary}
\begin{theorem}[Italiano--Martelli--Migliorini]
There exists a closed hyperbolic five-manifold $M^5$ admitting a smooth fibration $M^5\to S^1$ whose fiber $F^4$ is a closed aspherical four-manifold with $\pi_1(F^4)=H$ and $\pi_1(M^5)=G$ such that:
\begin{enumerate}[label=\textup{(\arabic*)}]
  \item $H\trianglelefteq G$ is a normal subgroup of the word-hyperbolic group $G$;
  \item $F^4$ is a finite $K(H,1)$ of geometric dimension~$4$;
  \item $H$ contains no Baumslag--Solitar subgroup $\mathrm{BS}(m,n)$ (since $G$ is hyperbolic and $H\le G$);
  \item $H$ is not word-hyperbolic (its outer automorphism action in $G/H\cong\mathbb{Z}$ has infinite order while $\mathrm{Out}(H)$ would be finite if $H$ were a non-elementary hyperbolic group with finite $K(H,1)$ of this type).
\end{enumerate}
\end{theorem}
\begin{remark}
Because $\dim F^4=4$, the strengthened variant of Question~1.1 restricted to groups of geometric dimension at most~$3$ remains open.
\end{remark}

\section{Methodological Framework and Verification Invariants}

The mathematical analysis for Result~1100101 in the domain of geometric group theory follows a strict separation between the exact statement decided and nearby formulations with stronger hypotheses. Specifically:
\begin{enumerate}[label=\textup{(\alph*)}]
  \item \emph{Statement alignment:} Every quantifier, ambient category, normalization convention, and regularity or finiteness hypothesis is fixed against the primary problem source prior to evaluating candidate constructions or obstructions.
  \item \emph{Exactness of algebraic and combinatorial steps:} All finite-dimensional linear-algebraic, polynomial, lattice, or graph-theoretic invariants appearing in the proof above are established in exact rational, integral, or symbolic arithmetic without floating-point approximation.
  \item \emph{Classification:} Within the program taxonomy, this result is categorized as a \emph{published counterexample} deciding the scope: dimension-unrestricted finite-classifying-space implications.
\end{enumerate}

\section{Scope and Open Boundaries}

To prevent over-claiming beyond the proved theorem, we record the explicit boundaries that are \emph{not} decided by the argument above:
\begin{itemize}
  \item finite classifying spaces of dimension at most three.
\end{itemize}

\begin{thebibliography}{9}
\bibitem{PrimarySource}
Primary problem list and source literature for Result~1100101 (geometric group theory), archived with exact theorem locators and verification certificates in the accompanying repository package.
\end{thebibliography}

\end{document}
